%=============================================================================!
% 13.7  WHICH THERMOSTAT, AND WHEN                                            !
%=============================================================================!
\section{Which thermostat, and when}
\label{sec:th-compare}

The sections so far measured each thermostat on its own. This section
sets the measurements side by side and draws from them which thermostat
to use for the three jobs a thermostat does: bringing a system to a
temperature, sampling averages at it, and leaving its dynamics fit to
measure.

\subsection*{Side by side}

\Cref{tab:th-compare} collects, for every thermostat and coupling, the
spread of $K$ in liquid argon as a fraction of the canonical spread, the
diffusion coefficient as a fraction of its value at fixed energy, and,
for one lithium atom with its own thermostat, how far the density of its
potential energy misses the exact one and how often it hops. The
diffusion coefficients carry the standard error of three runs, the
standard deviation of their three estimates divided by $\sqrt3$, since the variances of
the three add and their mean divides the sum by three; with so few runs it
is itself rough, and Chapter~15 gives the tools to say when two such
numbers differ. Three groups stand out.
\begin{itemize}
\item Rescaling fixes the kinetic temperature at each rescaling, and Berendsen
   coupling keeps its mean close to the target. Both narrow the spread, 0.11 to 0.59 of the canonical, and a single atom
   under Berendsen coupling stops hopping.
\item Andersen and Langevin sample the canonical ensemble once the
   coupling is \SI{1}{\per\pico\second} or more, for the liquid and for
   one atom alike, and cost no more than a step without a thermostat. At
   that coupling they cut the diffusion coefficient to 0.71 to 0.72 of its
   value, and at \SI{10}{\per\pico\second} to 0.17.
\item Nosé-Hoover chains and CSVR give the canonical spread at both
   couplings and sample one atom correctly. Like Berendsen coupling and a
   single Nosé-Hoover thermostat, they leave the diffusion coefficient
   1.01 to 1.07 of its value at fixed energy, in runs 0.4 to
   \SI{1.8}{\kelvin} warmer than those at fixed energy. The variances of
   independent quantities add, so the standard error of a difference is
   the square root of the sum of the two squared standard errors; the
   chains and CSVR lie at most 2.5 of these from fixed energy, and
   Berendsen coupling at \SI{1}{\pico\second} 2.9. Whether so small a
   difference is real is a question for Chapter~15. A single Nosé-Hoover thermostat matches the chains for
   the liquid but not for one atom. These findings agree with Basconi and
   Shirts, who compared thermostats on the transport properties of
   liquids\cite{Basconi2013}.
\end{itemize}

\begin{table}[htb]
\centering
\caption{The thermostats measured: for liquid argon at \SI{135}{\kelvin},
the relative spread of $K$ divided by the canonical
$\sqrt{2/N_{\mathrm f}}$, and
the diffusion coefficient divided by its value at fixed energy,
$(3.81\pm0.08)\times10^{-4}$\,\AA$^2$/fs (three runs of
\SI{100}{\pico\second}; rescaling, one); for one lithium atom on the
surface at \SI{1000}{\kelvin}, the largest gap between the density of its
potential energy and the exact one, divided by the exact density's peak,
and its hops per picosecond (1000 atoms over \SI{100}{\pico\second}).
For Andersen and Langevin, which move the centre of mass, $N_{\mathrm f} =
3N = 768$. Under Berendsen coupling the density is a narrow spike, so its
gap is many times the peak.}
\label{tab:th-compare}
\small
\begin{tabular}{@{}lllll@{}}
\toprule
& \multicolumn{2}{l}{liquid argon} & \multicolumn{2}{l}{one lithium atom}\\
thermostat & spread & $D$ & gap & hops/ps\\
\midrule
fixed energy & 0.61 & 1 & & \\
rescaling & 0.11 & 1.01 & & \\
Berendsen, \SI{0.1}{\pico\second} & 0.45 & $1.04\pm0.01$ & 15.5 & 0.001\\
Berendsen, \SI{1}{\pico\second} & 0.59 & $1.07\pm0.01$ & & \\
Andersen, \SI{0.1}{\per\pico\second} & 0.85 & $0.96\pm0.03$ & & \\
Andersen, \SI{1}{\per\pico\second} & 0.98 & $0.72\pm0.01$ & 0.011 & 0.974\\
Andersen, \SI{10}{\per\pico\second} & 1.00 & $0.173\pm0.004$ & & \\
Langevin, \SI{0.1}{\per\pico\second} & 0.93 & $0.97\pm0.02$ & 0.024 & 1.024\\
Langevin, \SI{1}{\per\pico\second} & 1.01 & $0.71\pm0.01$ & 0.005 & 1.008\\
Langevin, \SI{10}{\per\pico\second} & 1.01 & $0.169\pm0.001$ & 0.008 & 0.928\\
Nosé-Hoover, \SI{0.1}{\pico\second} & 0.99 & $1.02\pm0.01$ & 0.231 & 0.667\\
Nosé-Hoover, \SI{1}{\pico\second} & 1.04 & $1.05\pm0.01$ & & \\
chain of 3, \SI{0.1}{\pico\second} & 1.01 & $1.04\pm0.03$ & 0.016 & 1.087\\
chain of 3, \SI{1}{\pico\second} & 1.01 & $1.01\pm0.02$ & & \\
CSVR, \SI{0.1}{\pico\second} & 0.98 & $1.07\pm0.02$ & 0.012 & 1.081\\
CSVR, \SI{1}{\pico\second} & 0.99 & $1.05\pm0.01$ & & \\
\bottomrule
\end{tabular}
\end{table}

\subsection*{What to use}

\Cref{tab:th-choose} turns the measurements into a choice for each job.

\emph{Bringing a system to a temperature.} Any thermostat does it; what
matters is how fast it gets there and that it settles at the target. From
\SI{200}{\kelvin}, Berendsen and CSVR with $\tau_{\mathrm T} =
\SI{1}{\pico\second}$ and Langevin with $\gamma = \SI{1}{\per\pico\second}$
all reach the target within a few picoseconds (\cref{fig:th-berendsen}(b)).
Under Berendsen and CSVR the excess decays with the time constant
$\tau_{\mathrm T}C_V/C_K = \SI{1.59}{\pico\second}$, fitted at 1.67 and
\SI{1.63}{\pico\second}; under Langevin with $C_V/(2\gamma C_K) =
\SI{0.80}{\pico\second}$, fitted at \SI{0.68}{\pico\second}
(\cref{ex:th-langevin-tau}). A short coupling time, 0.1 to
\SI{1}{\pico\second}, is right here, and the run's early part is
discarded, for which Chapter~15 gives a measured criterion.

\emph{Averages at a temperature.} The average and the fluctuations must
both be canonical, so rescaling and Berendsen coupling are out. CSVR,
Nosé-Hoover chains, Langevin and Andersen all qualify. A single
Nosé-Hoover thermostat failed for one oscillator and for one atom on the
surface, and is expected to fail likewise for other small or nearly
harmonic systems, a single molecule or a nearly harmonic
crystal\cite{Martyna1992}. The heat capacity from the energy's
fluctuations, \cref{eq:sm-fluctuation}, shows what is at stake: it comes
out at 2.28 to $2.40N\kB$ under CSVR, the chain and Langevin, against the
$2.376N\kB$ of Chapter~12 from $\dd E/\dd T$, and at $0.41N\kB$ under
Berendsen coupling with $\tau_{\mathrm T} = \SI{0.1}{\pico\second}$.

\emph{Dynamics.} The thermostat must disturb the motion as little as
possible; of the motion, this chapter measured the diffusion of the
liquid. A run at fixed energy started from an equilibrated state disturbs
nothing, but its temperature is the one its energy gives, which in our
runs fell \SI{0.66}{\kelvin} short. CSVR or a Nosé-Hoover chain changes
the diffusion of the liquid by at most 7\%, and a coupling time of a
picosecond or more keeps the thermostat's hand light. Berendsen coupling
changes it no more, but it narrows the fluctuations, feeds the flying ice
cube and pins a single atom. Langevin and Andersen thermostats at
\SI{1}{\per\pico\second} remove more than a quarter of it.

\emph{An atom whose only bath is the thermostat.} When the real bath is
left out, as for an atom on a frozen surface, the thermostat's coupling
becomes part of the physics: the lithium atom's hop rate fell from 1.0 to
\SI{0.35}{\per\pico\second} as the friction rose from 3 to
\SI{100}{\per\pico\second}. Such a rate means something only
if the friction stands for a real bath, and the safer course is to
simulate the bath's atoms.

\begin{table}[htb]
\centering
\caption{Which thermostat to use, from the measurements of this chapter.}
\label{tab:th-choose}
\small
\begin{tabular}{@{}>{\raggedright\arraybackslash}p{0.25\linewidth}
   >{\raggedright\arraybackslash}p{0.38\linewidth}
   >{\raggedright\arraybackslash}p{0.3\linewidth}@{}}
\toprule
job & use & avoid\\
\midrule
bringing to a temperature & Berendsen, CSVR or Langevin,
   $\tau_{\mathrm T}$ of 0.1 to \SI{1}{\pico\second} or $\gamma$ of 1 to
   \SI{10}{\per\pico\second} (measured at \SI{1}{\pico\second} and
   \SI{1}{\per\pico\second}), the early part discarded & one Berendsen
   thermostat on a single atom, which pins its energy\\
averages and fluctuations & CSVR, Nosé-Hoover chains, Langevin, Andersen
   & rescaling, Berendsen; one Nosé-Hoover thermostat on small or nearly
   harmonic systems\\
dynamics (here, diffusion) & fixed energy after equilibrating; CSVR or a
   chain with $\tau_{\mathrm T} \ge \SI{1}{\pico\second}$ & Langevin and
   Andersen at $\ge$ \SI{1}{\per\pico\second}; Berendsen, for its
   fluctuations and flying ice cube\\
learned runs (TrajCast) & CSVR once per stride; its update preserves the canonical
   distribution of $K$ & Berendsen\\
\bottomrule
\end{tabular}
\end{table}

\begin{keyidea}
Rescaling and Berendsen coupling are for reaching a temperature. CSVR,
Nosé-Hoover chains, Langevin and Andersen sample the canonical ensemble; of
these,
CSVR and chains change the diffusion of the liquid by at most 7\%, while
friction or collisions at \SI{1}{\per\pico\second} slow it by more than
a quarter. A single Nosé-Hoover thermostat fails for small systems, and
for an atom whose bath is the thermostat, the atom's rates belong to the
coupling.
\end{keyidea}

\begin{selfcheck}
A run is to measure the diffusion of lithium in a liquid electrolyte at
\SI{300}{\kelvin}. Which thermostat, with what coupling, does this chapter
recommend for the production run, and why not Langevin with $\gamma =
\SI{5}{\per\pico\second}$?
\end{selfcheck}
